\section{Setup}\label{sec:setup}

\paragraph{PARA in brief.}
PARA \citep{kang2026para} combines two predictors for an image $x$ and classes $c = 1, \dots, C$.
The \emph{anchor} is the frozen zero-shot prediction
$z^0_c(x) = \tau_0 \langle F^0_v(x), t_c \rangle$, where $F^0_v$ is the frozen image encoder,
$t_c$ is the normalized mean embedding of an ensemble of class prompts, and $\tau_0$ is the frozen pretrained logit scale.
The \emph{support branch} is trained on the labelled support set $S$:
$z^S_c(x) = \tau_v \langle F^S_v(x), p^S_c \rangle$, with a LoRA-adapted encoder $F^S_v$ \citep{hu2022lora},
class prototypes $p^S_c$ and a learned scale $\tau_v$, trained with a leave-one-out prototype loss.
Each branch is standardized across classes (written $\bar z$), and the prediction moves from the anchor towards the support branch by a step $\eta(S)$:
\begin{equation}
  \hat z_c(x) = \bar z^0_c(x) + \eta(S)\,\big(\bar z^S_c(x) - \bar z^0_c(x)\big).
  \label{eq:blend}
\end{equation}
The step is computed in closed form from the frozen support embeddings $u_i = F^0_v(x_i)$, which, like $t_c$, are L2-normalized.
With $S_c$ the $n_c$ support images of class $c$, $h_c = \sum_{i \in S_c} u_i$ and $d_c = h_c / \lVert h_c \rVert$,
\begin{equation}
  \eta(S) = \frac{\lambda}{C} \sum_{c=1}^{C} \frac{\delta_c \phi_c}{1 + \delta_c \phi_c},
  \qquad
  \delta_c = 1 - \langle d_c, t_c \rangle,
  \qquad
  \phi_c = \frac{\lVert h_c \rVert^2}{n_c}.
  \label{eq:eta}
\end{equation}
The \emph{anchor discrepancy} $\delta_c$ is meant to measure how far the class images lie from their own text,
and the \emph{support agreement} $\phi_c$ how consistently they point the same way.
$\lambda = 0.9$ on every dataset.

\paragraph{Reimplementation.}
No official code has been released, so we reimplemented PARA from the paper.
The backbone is BiomedCLIP ViT-B/16 \citep{zhang2023biomedclip}, with $\tau_0 = 85.23$.
LoRA uses rank $8$ and $\alpha = 16$ on the image encoder.
Following the paper, we train for $80$ epochs with AdamW (learning rate $10^{-4}$ for LoRA and $10^{-3}$ for $\tau_v$, weight decay $10^{-4}$)
and keep the final checkpoint.
Two details the paper leaves open we set ourselves: a batch size of $32$ and random resized crops with scale $(0.5, 1)$ plus horizontal flips.
As a check, BiomedCoOp's GPT-4 prompts give a DermaMNIST zero-shot accuracy of $53.55\%$ in our implementation,
against $53.78\%$ for zero-shot BiomedCLIP in Table A3 of \citet{kang2026para}.
This confirms that the pipeline matches theirs, not that the anchor is useful: both are below the $66.7\%$ obtained by always predicting
the majority class (nevi), and the corresponding BAcc is $21.39\%$ (Table~\ref{tab:T1_datasets}).

\paragraph{Datasets and prompts.}
We use six 2D datasets from MedMNIST v2 \citep{yang2023medmnist} (Table~\ref{tab:T1_datasets}):
DermaMNIST at $224$ pixels and the other five at $128$ pixels (all resized to the encoder's $224$-pixel input).
Each dataset is scored on a fixed test subset.
Because splits and resolution differ from the paper's benchmark, our absolute numbers are not directly comparable to it.
Every dataset uses our own prompt ensemble:
five hand-written visual descriptions per class for DermaMNIST, and five generic templates for the other datasets.
Where BiomedCoOp \citep{koleilat2025biomedcoop} provides GPT-4 descriptions for exactly matching classes,
we add them as a second prompt set (50 per class): DermaMNIST, OCTMNIST and BloodMNIST,
and BreastMNIST via its source dataset BUSI, with BUSI's ``benign'' and ``normal'' descriptions pooled into MedMNIST's merged class.

\begin{table}[t]
\centering
\small
\caption{Datasets. Zero-shot BAcc (\%) on the evaluated test subset, per prompt set. ``Your prompts'' are 5 hand-written visual descriptions per class for DermaMNIST and 5 generic templates elsewhere; GPT-4 prompts are BiomedCoOp's (BreastMNIST via BUSI: benign and normal pooled). ZS: zero-shot.}
\label{tab:T1_datasets}
\begin{adjustbox}{max width=\linewidth}
\begin{tabular}{lrrrrrl}
\toprule
Dataset & Classes & Test images & Pixels & ZS BAcc, own & ZS BAcc, GPT-4 & Split \\
\midrule
DermaMNIST & 7 & 2000 & 224 & 20.93 & 21.39 & design \\
OCTMNIST & 4 & 1000 & 128 & 49.30 & 51.50 & design \\
BloodMNIST & 8 & 2000 & 128 & 14.86 & 13.49 & design \\
BreastMNIST & 2 & 156 & 128 & 53.63 & 70.93 & held-out \\
PneumoniaMNIST & 2 & 624 & 128 & 67.35 & -- & held-out \\
OrganAMNIST & 11 & 2000 & 128 & 25.71 & -- & held-out \\
\bottomrule
\end{tabular}
\end{adjustbox}
\end{table}


\paragraph{Protocol.}
For each dataset we use $K \in \{1, 4, 16\}$ shots per class and draw 10 support sets per $K$.
The draws are made once and shared by every step rule and both prompt sets, and the support branch is trained once per draw
(one training seed) and cached.
Every step rule below is therefore applied to identical trained models, and differences between rules are paired by draw.
Our primary metric is balanced accuracy (BAcc), because plain accuracy rewards majority-class predictions
(on DermaMNIST, about two thirds of test images are nevi).
We also report CVaR20, the mean BAcc over the worst $20\%$ of draws (the worst 2 of 10), which is PARA's own worst-case measure,
and Macro-F1.
For paired comparisons we run a Wilcoxon signed-rank test over the 10 draws of each setting (dataset, prompt set, $K$),
count settings significantly better or worse at $p < 0.05$, and also report counts after Holm's correction within each family of settings.
With 10 draws the smallest attainable two-sided $p$ is $2/1024 \approx 0.002$, so Holm-significant settings are those
in which nearly every draw moves in the same direction.

\paragraph{Step rules compared.}
Besides PARA's $\eta(S)$ we evaluate, on the same trained models:
zero-shot ($\eta = 0$); visual only ($\eta = 1$);
a \emph{shot-only} schedule $\eta_K = \lambda K / (1 + K)$, which uses the shot count and nothing else
(shot-dependent schedules are prior art, \eg ProLIP \citep{fahes2026prolip}, so we treat it as a baseline);
a HOSO-style step chosen on one held-out shot per class \citep{vorster2026hoso};
and an oracle that picks the best $\eta$ on the test set for each draw, reported only as an upper bound.

\paragraph{Design and held-out datasets.}
DermaMNIST, OCTMNIST and BloodMNIST were used to develop every analysis and rule (\emph{design} datasets).
PneumoniaMNIST, BreastMNIST and OrganAMNIST were chosen afterwards from a zero-shot screening alone,
before any model was trained on them, and serve as \emph{held-out} datasets.
The one tuned rule in this paper, the anchor-aware step of Section~\ref{sec:finding3},
has a single constant $c$. We fixed $c = 8$ because it maximized CVaR20 BAcc pooled over the three design datasets
($56.07$, against $55.71$ for $c = 4$ and $55.61$ for $c = 16$); this choice is in-sample on those datasets.
We made this choice between the last design analysis and the first held-out training run, but the evidence for that timing is weak.
The decision itself was not recorded with a date: the held-out runs scored all nine values of $c$ without marking one as chosen,
and the first committed record of $c = 8$ postdates the held-out runs.
The zero-shot screening, which also computed $r_a$ on the held-out candidates, may have preceded the choice.
We therefore call it a pre-specified rule, not a registered one; that $c = 8$ is not the best value on the held-out data
(Section~\ref{sec:finding3}) argues against a choice informed by held-out results, but does not prove the order.
